\section{Discussion}

\emph{The answer.} Eq.~3, Eq.~4, and the very sparse matrices of Li, Hastie, and Church at these small $d$ show no loss against a Gaussian that survives correction, except Eq.~4 on four-feature Iris, where sparsity drops informative features and can leave a column of the matrix all zero (the loss comes from the draws with a zero column). Equivalence within $\pm 2$ points is shown mostly at the larger $k$. In paired differences pooled over draws, target sizes, and the data sets other than Iris on which each family was run, they trail by under a point (Eq.~4 by \N{mag:pool:eq4:gauss} with Iris), which is why the statement is ``within about a point of dense'' outside Iris. One nonzero per column trails at every target size.

\emph{Why the original advantage appeared.} The report compared two matrices from one draw, on one split. About half of all draws put Eq.~4 ahead on three or more data sets, but over 100 draws, the mean difference has no consistent sign. Separately, choosing kernels on the test set, as the report did, makes accuracies look better: in the tuned study, it adds about \N{z:test_optimism} points of plain accuracy on average. It applies to both matrices, so it is mainly a caveat about the level of the accuracies, not an explanation of the gap between them.

\emph{Distortion.} Theory gave no reason to expect Eq.~4 to beat Eq.~3 or the Gaussian, since Eq.~4's $q$ has the Gaussian's variance. It also gave no reason to expect it to lose to the Gaussian. Apart from one nonzero per column, which has both the highest variance and the lowest accuracy, the variance order does not explain the accuracy order, and ``distortion explains the ordering'' is not supported.
